\documentclass[10pt]{article}

\usepackage[
  paperwidth=160mm,
  paperheight=90mm,
  margin=0mm
]{geometry}
\usepackage{amsmath,amssymb}
\usepackage[T1]{fontenc}
\usepackage{lmodern}
\usepackage{microtype}
\usepackage{xcolor}
\usepackage{tikz}
\usepackage[hidelinks]{hyperref}

\definecolor{ink}{HTML}{14213D}
\definecolor{blue}{HTML}{2864DC}
\definecolor{bluepale}{HTML}{EEF4FF}
\definecolor{muted}{HTML}{52627A}
\definecolor{paper}{HTML}{FAFBFE}

\pagestyle{empty}
\setlength{\parindent}{0pt}
\renewcommand{\familydefault}{\sfdefault}

\begin{document}
\color{ink}

\begin{tikzpicture}[remember picture,overlay]
  \fill[paper] (current page.south west) rectangle (current page.north east);
  \fill[blue] (current page.north west) rectangle
    ([yshift=-2.2mm]current page.north east);

  \node[anchor=north west,inner sep=0pt,text width=146mm]
    at ([xshift=7mm,yshift=-7mm]current page.north west) {%
      {\bfseries\fontsize{17.28}{19}\selectfont
      3SAT $\longrightarrow$ zero-finding with constant Jacobian}
    };

  \node[anchor=north west,inner sep=0pt,text width=146mm]
    at ([xshift=7mm,yshift=-15mm]current page.north west) {%
      {\color{muted}\fontsize{6.2}{7}\selectfont
      Source gadget: Alpöge's announced Jacobian-conjecture counterexample
      · \href{https://x.com/__alpoge__/status/2079028340955197566}%
      {x.com/\_\_alpoge\_\_/status/2079028340955197566}}
    };

  \node[anchor=north west,fill=bluepale,inner sep=1.7mm,text width=142.6mm]
    at ([xshift=7mm,yshift=-20mm]current page.north west) {%
      {\fontsize{8.5}{9.7}\selectfont
      \textcolor{blue}{\textbf{CLAIM.}}
      Given a $3$CNF formula $\phi$, construct a degree-at-most-$7$ map
      $K_\phi$ with
      \setlength{\abovedisplayskip}{0.5mm}
      \setlength{\belowdisplayskip}{0mm}
      \[
        K_\phi^{-1}(0)\ne\varnothing
        \iff \phi\in\mathrm{SAT},
        \qquad
        \det JK_\phi\equiv(-3456)^{n+m}.
      \]}
    };

  \node[anchor=north west,inner sep=0pt,text width=146mm]
    at ([xshift=7mm,yshift=-35mm]current page.north west) {%
      {\fontsize{7.2}{8.35}\selectfont
      \setlength{\abovedisplayskip}{0.7mm}
      \setlength{\belowdisplayskip}{0.7mm}
      \textbf{1. The gadget.}
      Fix $H:\mathbb R^3\to\mathbb R^3$ with $\det JH\equiv-3456$,
      $H^{-1}(0)=\{r_0,r_+,r_-\}$, and one omitted target
      $\Delta\notin\operatorname{im}(H)$.  The decoder $\beta(x,y,z)=x^2$
      reads $r_0$ as $0$ and $r_\pm$ as $1$.

      \vspace{0.5mm}
      \textbf{2. One clause and the switch.}
      Define the two-input map
      $T(v,s)=H(v)-(1-s)\Delta$.  The equation $T(v,s)=0$ asks $H$ to output
      $(1-s)\Delta$: $s=1$ asks for $0$ (possible), while $s=0$ asks for
      $\Delta$ (impossible).  For example, if
      $C_j=X_a\lor\neg X_b\lor X_c$, set
      \[
        s_j=1-(1-\beta(w_a))\,\beta(w_b)\,(1-\beta(w_c)).
      \]
      This is $1$ exactly when the clause is true, so its block satisfies
      \[
        \exists v_j:\ T(v_j,s_j)=0
        \quad\Longleftrightarrow\quad C_j\text{ is true}.
      \]

      \textbf{3. The stack.}
      Put one $H(w_i)$ block per variable and one $T(v_j,s_j(w))$ block per
      clause:
      \[
        K_\phi=
        \bigl(H(w_1),\ldots,H(w_n),
        T(v_1,s_1(w)),\ldots,T(v_m,s_m(w))\bigr).
      \]
      The $H$ blocks encode one assignment; the $T$ blocks vanish exactly when
      all its clauses are true.  Hence $K_\phi=0\iff\phi$ is satisfiable.
      Finally, $s$ only translates the output of $H$, so
      $D_vT(v,s)=JH(v)$.  The full Jacobian is block lower triangular with
      $n+m$ copies of $JH$ on the diagonal, giving
      $\det JK_\phi=(-3456)^{n+m}$ and $\deg K_\phi\le7$.}
    };
\end{tikzpicture}

\null
\end{document}
